\documentclass[12pt, a4paper]{article}
\usepackage[utf8]{inputenc}
\usepackage[T1]{fontenc}
\usepackage{geometry}
\geometry{a4paper, margin=2.5cm}
\usepackage{amsmath, amssymb}
\usepackage{graphicx}
\usepackage{booktabs}
\usepackage{hyperref}
\usepackage{setspace}
\usepackage{titlesec}
\usepackage{lmodern}

\onehalfspacing
\titleformat{\section}{\large\bfseries}{\thesection}{1em}{}
\titleformat{\subsection}{\normalsize\bfseries}{\thesubsection}{1em}{}

\title{Formal Comparison of Critical Behavior: \\ Plausible (GSS) vs. Random (ER) Topologies}
\author{A Theoretical Evaluation of Topological Destruction}
\date{\today}

\begin{document}

\maketitle

\begin{abstract}
This report provides a formal mathematical and physical evaluation of the phase transition behaviors observed in a social contagion model simulated over empirically grounded plausible networks (GSS) versus appropriately randomized variants (Erdős-Rényi, ER). By analyzing the critical susceptibility exponents ($\gamma$) and the macroscopic order parameter macroscopic jumps ($\Delta \Phi$), we demonstrate that the destruction of homophilic structural clustering results not in a partial degrading of network diffusion capacity, but in a complete thermodynamic collapse of the universality class. We establish that plausible networks exhibit continuous (second-order) phase transitions belonging to the Mean-Field Ising universality class ($\gamma \approx 1$), while random networks collapse into trivial, discontinuous (first-order) cascades.
\end{abstract}

\section{Introduction and Thermodynamic Framework}

The study of collective social contagion can be mapped onto the formalisms of statistical physics by defining an order parameter $\Phi \in [0, 1]$, representing the global proportion of adopters, and an external control parameter $MSP$ (Maximum Social Proximity), acting as the systemic thermal driving force (an inverse temperature analogue) against the intrinsic rigidity defined by individual thresholds (IUL). 

In the thermodynamic limit ($N \to \infty$), the system's susceptibility $\chi$, which measures the responsiveness of the order parameter to infinitesimal shifts in the control parameter, is governed by the Fluctuation-Dissipation Theorem (FDT):
\begin{equation}
    \chi \propto N \cdot \text{Var}(\Phi)
\end{equation}

Near a critical boundary $MSP_c$, if the system undergoes a continuous (second-order) phase transition, the susceptibility diverges according to a universal power law:
\begin{equation}
    \chi \sim |MSP - MSP_c|^{-\gamma}
\end{equation}

The critical exponent $\gamma$ exclusively defines the ``Universality Class'' of the system. Its value depends strictly on the fundamental symmetries and the spatial dimensionality of the interactions, ignoring the micro-level nuances of the specific model.

\section{Plausible Networks (GSS): The Richness of Continuous Criticality}

For the plausible geometries imputed holding the empirical General Social Survey (GSS) parameters, diffusion exhibits a macroscopic continuous phase shift. Over the non-biased analogue slice (IUL $\approx 0.200$), the system yields an exponent of:
\begin{equation}
    \gamma_{GSS} \approx 0.977 \pm 0.01
\end{equation}

This value aligns extraordinarily well with the theoretical expected value for the \textbf{Mean-Field Ising Universality Class} ($\gamma_{MF} \equiv 1.0$). 

\subsection{The Structural Meaning of $\gamma \approx 1$}
The mathematical richness of possessing a $\gamma \approx 1.0$ is immense. It proves that the human homophilic topology is capable of sustaining \textit{scale-free} fractal fluctuations. In the critical zone, the network contains highly connected homophilic ``echo chambers'' (clusters of localized resistance or facilitation) that act as macroscopic domain walls. 

These domain walls impose a structural friction that forces the order parameter $\Phi$ to grow continuously. Avalanches of adoption occur at all possible size scales (power-law distribution of cascades), allowing the variance to act smoothly as a predictive precursor to systemic phase change. The universality class validates that societal coordination is not an algorithmic artifact, but a fundamental emergent property of complex systems behaving as a unified whole.

\section{Topological Destruction (GSS-ER) and the Loss of Universality}

To evaluate the pure causal effect of the empirical structure, the GSS topologies were deliberately destroyed via random rewiring, preserving only the macroscopic density whilst imposing a Poissonian degree distribution (Erdős-Rényi, ER).

The destruction of the structure is not partial; it is \textbf{total} from a thermodynamic perspective. For the GSS-ER control networks, the susceptibility exponent inflates erratically:
\begin{equation}
    \gamma_{ER} \approx 3.7 \text{ to } 4.6
\end{equation}

These values do not belong to any formal Universality Class in statistical physics. More critically, for the super-critical regime ($MSP > MSP_c$), the variance of the order parameter collapses to mathematically exactly zero ($\text{Var}(\Phi_{ER}) = 0$).

\subsection{Collapse into First-Order Discontinuity}
The ER results prove the absence of a continuous phase shift. By destroying the local clustering, the ER network drastically reduces the average path length and eliminates the echo-chamber domain walls. 

Consequently, once the energetic boundary $MSP_c$ is marginally breached, the random network cannot offer any local friction capable of stopping the contagion. Every identically initialized Monte Carlo seed triggers an immediate, unmitigated global cascade ($\Phi \to 1$). This describes a \textbf{First-Order (Discontinuous) Phase Transition}. Because the jump is instantaneous, there are no scale-free fractal fluctuations, no gradual continuous divergences, and the susceptibility drops to zero post-jump. The calculated $\gamma_{ER}$ exponents are simply computational artifacts of fitting an invalid continuous regression to a fundamentally discontinuous wall.

\section{Conclusion: What is Lost?}

The comparative analysis reveals that substituting an empirically plausible network with a random topology destroys much more than localized demographic realism; it extinguishes the entire thermodynamic character of the social system. 

The richness lost is the \textbf{complexity} itself. Random networks render the population as a mathematically trivial conduit where information propagates identically and predictably. In stark contrast, the $\gamma \approx 1.0$ derived from the GSS matrices demonstrates that empirical social topologies uniquely construct a continuous structural tension between localized identity preservation (homophily) and systemic coordination. The continuous criticality relies \textit{entirely} on the non-random arrangement of societal ties. 

Therefore, we conclude that the capacity of a society to undergo rich, non-linear, self-organized behavioral transitions is exclusively scaffolded by its topological geometry. The network, quite literally, decides the fundamental laws of physics the societal collective will obey.

\end{document}
